  
\addcontentsline{toc}{chapter}{\twkkai{对话录·其四}}  
\markboth{\twkkai{对话录·其四}}{}  
\pagecolor{black}
\color{white}
\quad\quad \twkkai{\lettrine[,lraise=0.2,lhang=0.5]{\textbf{猫}}{：喵！喵！}

狄利克雷：喂！回来！你过头了！可别太过分！

猫：喵……呼哧，呼哧……

狄利克雷：何必呢？搞得自己这么累……不要让自己变成人间的罪人；理智点。

猫：好啦；在咖啡馆听到一点激昂的音乐，肾上腺素有点过量了……

狄利克雷：我看你只是厌烦了和我的生活，不捣点乱都不行……

猫：我知错了，我的朋友。啊，你看到人间的那炫目的音乐喷泉了吗？

狄利克雷：我可什么都看不到呢。你又有什么奇妙见闻啦？

猫：我猜，正是那喷泉和它那塞壬之乐——韦伯Op.73第一乐章——导致了我刚才的疯狂行径。

狄利克雷：你又在找借口了。正事不做，脑子里的想法倒是一大堆……

（猫哼着小调，踮着脚悄悄地跳起了舞\marginnote{}{\small ♪}。）

\rule{0pt}{5pt}

\noindent\hspace*{3cm}\includegraphics[height=11mm]{frontmatter/猫.pdf}

狄利克雷：这音乐果然有如此神力吗——能让你刚才行为失常……

猫：当然。尤其是低音提琴的这一句\kern-15pt\lily[hpadding=0pt]{
    \override Staff.StaffSymbol.color       = #white
    \override Staff.Clef.color              = #white
    \override Staff.TimeSignature.color     = #white
    \override Staff.KeySignature.color      = #white
    \override Staff.BarLine.color           = #white
    \override Staff.LedgerLineSpanner.color = #white
    \override NoteHead.color                = #white
    \override Stem.color                    = #white
    \override Beam.color                    = #white
    \override Flag.color                    = #white
    \override Rest.color                    = #white
    \override Accidental.color              = #white
    \override Dots.color                    = #white
    \override Slur.color                    = #white
    \override Tie.color                     = #white
    \override Script.color                  = #white
    \override DynamicText.color             = #white
    \override Hairpin.color                 = #white
    \override TextScript.color              = #white
    \override TupletBracket.color           = #white
    \override TupletNumber.color            = #white
  \clef bass
  \key as \major
  \time 3/4
  f,2 as,8. as,16 |
  c8 r8 as,8 r8 f,8 r8 |
  e,4 -. c4 -. r4 |
  des4 ( -> c4 ) -! r4
  \bar "||"
}\kern5pt 真的让我从一开始就心潮澎湃了！

狄利克雷：的确，从听感上就很有韦伯的风格——不拘于古典。

猫：你应该想说的是……巴洛克风格吧？

狄利克雷：是的；看来你造诣深一点。

猫：怎敢怎敢！

狄利克雷：回到正题吧。巴洛克风格……之后呢？

猫：喵——这句利都奈洛主题\index{lidunailuo@利都奈洛主题（Ritornello）}\kern-15pt\lily[hpadding=0pt]{
    \override Staff.StaffSymbol.color       = #white
    \override Staff.Clef.color              = #white
    \override Staff.TimeSignature.color     = #white
    \override Staff.KeySignature.color      = #white
    \override Staff.BarLine.color           = #white
    \override NoteHead.color                = #white
    \override Stem.color                    = #white
    \override Beam.color                    = #white
    \override Flag.color                    = #white
    \override Rest.color                    = #white
    \override Accidental.color              = #white
    \override Dots.color                    = #white
    \override Slur.color                    = #white
    \override Tie.color                     = #white
    \override Script.color                  = #white
    \override DynamicText.color             = #white
    \override Hairpin.color                 = #white
    \override TextScript.color              = #white
    \override TupletBracket.color           = #white
    \override TupletNumber.color            = #white
  \clef bass
  \key as \major
  \time 3/4
  f,2 as,8. as,16 |
  c8 r8 as,8 r8 f,8 r8 |
  e,4 -. c4 -. r4 |
  des4 ( -> c4 ) -! r4
  \bar "||"
}\kern5pt （好吧，原谅我再重复了一次），韦伯这位伟大的作曲家多次利用了它，以推动呈示部、发展部和再现部的进行。而且，这里有个很有意思的点。

狄利克雷：是什么？

\clearpage
\thisblankpage
\includepdf[
  pages={1-3},
  nup=1x3,
  scale=0.7,
  pagecommand={
    \begin{tikzpicture}[remember picture,overlay]
      \node[align=center]
      at ($(current page.north east)+(-1.5cm,-12cm)$)
      {
        {\Large ♪}\\[0.8em]
        \qraudio
          {weber1-1}
          {\WeberSound}
          {https://raw.githubusercontent.com/kelvinlamresearchmath-web/finding-dirichlet-site/master/weber1-1.mp3}
      };
    \end{tikzpicture}
  }
]{weber_all.pdf}
\thisblankpage
\includepdf[
  pages={4-6},
  nup=1x3,
  scale=0.7
]{weber_all.pdf}

猫：让我揭晓吧——是韦伯在再现部违背了传统，让第二主题缺席了！

狄利克雷：第二主题？你哼一下。

猫：好的——大概就是\kern-15pt\lily[hpadding=0pt]{
  \relative bes'' {
    \override Staff.StaffSymbol.color       = #white
    \override Staff.Clef.color              = #white
    \override Staff.TimeSignature.color     = #white
    \override Staff.KeySignature.color      = #white
    \override Staff.BarLine.color           = #white
    \override Staff.LedgerLineSpanner.color = #white
    \override NoteHead.color                = #white
    \override Stem.color                    = #white
    \override Beam.color                    = #white
    \override Flag.color                    = #white
    \override Rest.color                    = #white
    \override Accidental.color              = #white
    \override Dots.color                    = #white
    \override Slur.color                    = #white
    \override Tie.color                     = #white
    \override Script.color                  = #white
    \override DynamicText.color             = #white
    \override Hairpin.color                 = #white
    \override TextScript.color              = #white
    \override TupletBracket.color           = #white
    \override TupletNumber.color            = #white
    \clef "treble" \time 3/4 \key bes \major \transposition bes
    \acciaccatura bes,8 bes'2( a8 g8 |
    f4. e8 es8 c8 |
    bes4.. a16 bes8. c16) |
    d2( c8) r8 |
    \bar "||"
  }
}\kern5pt 。它在呈示部中的关系大调（降A大调）出现过一次；本来应该在再现部呼应回去，把它拉回主调的——结果韦伯直接忽略了这一重大仪式！这就违反了奏鸣曲\index{zoumingqushi@奏鸣曲式（Sonata form）}\index{zoumingqushi@奏鸣曲式（Sonata form）!呈示部、发展部、再现部}\index{zoumingqushi@奏鸣曲式（Sonata form）!韦伯再现部第二主题的缺席}式的传统逻辑，让它变得“巴洛克”了。

狄利克雷：原来如此。

守门人：……把所有包含这些指标的\(L\)-函数乘起来，得到一个新的、有序的整体——我们的拼图便拼好了。在这里，缺口就等同于有其中一个\(L\)-函数在\(s=1\)的时候等于0——利用的是反证法……

猫：这守门人——第几次梦呓啦？

狄利克雷：你倒不用管他；而且我觉得，像你之前所说，他每次都呼应着我们这里的音乐主题呢。

猫：倒是，我也觉察到了。而且讲的是你发现的算术级数中的素数定理\index{dilikelei@狄利克雷（P.\,G.\,L.\ Dirichlet）!算术级数中的素数定理}里最关键的一步呢。

狄利克雷：是的。我觉得你又有一大堆理论要讲了对不对？

猫：可真了解我。你看，利都奈洛主题就像是守门人所说的那个“新的、有序的整体”——\(\zeta_m(s)\)所起到的作用。                                

狄利克雷：怎么说？

猫：你看嘛——它最初的表达式\(\zeta_m(s)=\prod_{\chi\in\widehat{G}}L(s,\chi)\)是为了让所有\(L\)-函数拼在一起，通过整体拼图找出缺口。具体来说，就是先找到某个\(L\)-函数在\(s=1\)的时候为零。然后就可以利用反证法了——不过这是题外话。总之，利都奈洛主题也是这样——在每一次乐章段落的转换时都需要出现一次，以起到衔接整体乐章\index{lidunailuo@利都奈洛主题（ritornello）!衔接乐章段落的作用}的作用。而且这里还有个更妙的联系。

狄利克雷：说说看，不用等我的——毕竟是我先写出来的。

猫：没有没有，不是怀疑你的意思。我是想说，这条表达式\(\zeta_m(s)\\=\prod_{\chi\in\widehat{G}}\prod_{p\,\nmid\, m}\left(1-\frac{\chi(p)}{p^{s}}\right)^{-1}=\prod_{p\,\nmid\, m}\prod_{\chi\in\widehat{G}}\left(1-\frac{\chi(p)}{p^{s}}\right)^{-1}\)中交换\(\prod_\chi\)和\(\prod_p\)顺序的步骤，简直是精准对应了一句话——

守门人：……我们要让某种对称性自己浮出水面。

猫：这守门人，老抢我台词！不过他又说到点上了。的确，先遍历\(\chi\)再遍历\(p\)，和先遍历\(p\)再遍历\(\chi\)，两者不是一回事。前者很难被分析；后者有一种对称性，而且这样就可以对内层求积\(\prod_{\chi\in\widehat{G}}\left(1-\frac{\chi(p)}{p^{s}}\right)^{-1}\)进行分析了——\(\chi(p)\)取遍单位根的这一结论很快就呼之欲出了。这和那个利都奈洛主题也是密切相关的。

狄利克雷：你说说。

猫：（喘着气）利都奈洛主题在发展部出现，藏在织体里，和独奏声部并行，不以全奏形式强调自己——这不就是一种隐蔽和潜行吗？你需要换一种听觉思维才能听到（或者说，捕捉）到它！整体上听的东西没有变，但是你捕捉的重点变了！\index{lidunailuo@利都奈洛主题（ritornello）!在发展部的隐藏与对称性类比}

狄利克雷：嗯，你说得很清楚了。本质上就是：原本表达的结构不利于我们发现隐蔽的对称性，需要我们自己巧妙地重新排列，它才会自己显形。

猫：是的！就是指换一个站位，原本藏着的东西才会出现。而且！这或许跟使者2一直没有在天庭出现过息息相关。

狄利克雷：你这站位换得也太快了吧……？虽然使者的问题的确值得我们一起斟酌，但你的联想能力我也是不得不佩服——老喽，老喽！

猫：哎哟，您老就别说丧气话了。

狄利克雷：不过有个问题——为什么你刚才演奏的是第一乐章，之前先演奏第二乐章呢？

猫：好问题——而且我觉得你应该自己回答自己呢！你忘啦——你去过人间扰乱过时间嘛；也会扰乱天庭的。

狄利克雷：是这么回事吗？好——不过你演的真的不错。

猫：谢谢喵！而且，你可知道这个乐章后来启发了你大舅子门德尔松的创作喵？

狄利克雷：嗯……还有这样的事吗？看来我对他实在太不了解了——回去倒是要问问他。

猫：莫急！下次让我演奏完门德尔松你再找他也不迟。

狄利克雷：为何？

猫：免得你空手而去，不知道和他聊什么嘛。


\rule{0pt}{20pt}

\noindent\adjustbox{center}{\includegraphics[scale=1]{trebleclef.preview.pdf}}

}


